% !TEX root = hw5-answers-graded.tex
\chead{\textbf{Exercises week 5}}

\begin{document}
\flushleft\textbf{Topics: d-separation, $\sigma$-separation, acyclification}
\begin{enumerate}[label=5.\arabic*]
    \item %(1 + 1 points)
    Given a graph (where $J,K$ are input nodes)
    \[
        \begin{tikzpicture}[scale=1, transform shape]
            \node[ndout] (v1) at (5,5) {$A$};
            \node[ndout] (v2) at (3,3) {$B$};
            \node[ndout] (v3) at (7,3) {$D$};
            \node[ndout] (v7) at (11,3) {$E$};
            \node[ndout] (v5) at (5,1) {$F$};
            \node[ndout] (v6) at (7,-1) {$G$};
            \node[ndout] (v12) at (4,6) {$H$};
            \node[ndout] (v13) at (6,6) {$I$};
            \node[ndint] (v4) at (1,4) {$J$};
            \node[ndint] (v11) at (11,5) {$K$};
            \draw[arint] (v4) to (v2);
            \draw[arint] (v11) to (v7);
            \draw[arlat, bend left] (v1) to (v3);
            \draw[arlat, bend left] (v1) to (v7);
            \draw[arlat] (v2) to (v3);
            \draw[arout] (v1) to (v2);
            \draw[arout] (v1) to (v12);
            \draw[arout] (v12) to (v13);
            \draw[arout] (v13) to (v1);
            \draw[arout] (v1) to (v3);
            \draw[arout] (v2) to (v5);
            \draw[arout] (v3) to (v5);
            \draw[arout] (v5) to (v6);
            \draw[arout] (v7) to (v6);
        \end{tikzpicture}
    \]
    \begin{enumerate}
        \item We define the length of a walk/path as the number of edges in the walk/path. For this exercise, also consider walks/paths to be equal when they are in reverse order, e.g.\ $A \to B$ is equal to $B \leftarrow A$.
        \begin{enumerate}
            \item List all bidirected paths of length 3.
            \item Give two examples of bidirected walks of length 3.
            \item List all collider paths of length 5.
            \item Give two examples of a collider walk of length 5.
            % \item Given an example of a trek of length 4.
        \end{enumerate}
        \item Find for node $A$ the
        \begin{enumerate}
            \item parents
            \item children
            \item siblings
            \item ancestors
            \item descendants
            \item strongly connected component
            \item district.
        \end{enumerate}
        % \item Find for each node in the graph the (d-separation) Markov blanket. (1pt)
    \end{enumerate}

    \begin{ungradedsolution}
        \begin{enumerate}
            \item
            \begin{enumerate}
                \item Bidirected paths: $B\leftrightarrow D \leftrightarrow A \leftrightarrow E$
                \item Bidirected walks: $B\leftrightarrow D \leftrightarrow A \leftrightarrow E$ and $B\leftrightarrow D \leftrightarrow A \leftrightarrow D$
                \item Collider paths: $J \rightarrow B \leftrightarrow D \leftrightarrow A \leftrightarrow E \leftarrow K$
                \item Collider walk: $A \rightarrow B\leftrightarrow D \leftrightarrow A \leftrightarrow E \leftarrow K$ and $I \rightarrow A \leftrightarrow D \leftrightarrow A \leftrightarrow E \leftarrow K$
                % \item Trek: $F\leftarrow B\leftarrow A \leftrightarrow E \rightarrow G$.
            \end{enumerate}
            \item
            \begin{enumerate}
                \item $\Pa(A) = \{I\}$
                \item $\Ch(A) = \{H, B, D\}$
                \item $\Sib(A) = \{D, E\}$
                \item $\Anc(A) = \{H, I, A\}$
                \item $\Desc(A) = \{H, I, A, B, D, F, G\}$
                \item $\Sc(A) = \{H, I, A\}$
                \item $\Dist(A) = \{B, D, A, E\}$
            \end{enumerate}
        \end{enumerate}
    \end{ungradedsolution}

    \item %(1 + 1 points)
    Consider the following CDMG $G$, where $A,H$ are input nodes:
    \begin{figure}[ht]
        \centering
        \begin{tikzpicture}[scale=1, transform shape]
            \node[ndout] (B) at (2, 0) {$B$};
            \node[ndout] (C) at (4, 0) {$C$};
            \node[ndout] (D) at (6, 0) {$D$};
            \node[ndout] (E) at (8, 0) {$E$};
            \node[ndout] (F) at (10, 0) {$F$};
            \node[ndint] (A) at (0, 0) {$A$};
            \node[ndint] (H) at (8, -2) {$H$};
            \draw[arint] (A) to (B);
            \draw[arout] (C) to (B);
            \draw[arout] (C) to (D);
            \draw[arout] (E) to (D);
            \draw[arout] (E) to (F);
            \draw[arint] (H) to (E);
        \end{tikzpicture}
    \end{figure}
    % \begin{enumerate}
    % \item

    Do the following $d$-separations hold on graph $G$? Give a brief explanation for each answer.
    \emph{Note that we write $C \Perp^d E \mid D$ for $\{C\} \Perp^d \{E\} \mid \{D\}$.}
    \begin{enumerate}
        \item $\displaystyle  C \Perp^d E$
        \item $\displaystyle C \Perp^d E \mid D$
        \item $\displaystyle B \Perp^d D \mid \{A, H\}$
        \item $\displaystyle B \Perp^d D \mid \{C, A, H\}$
    \end{enumerate}
    % \item Does $d$-separation $\displaystyle B \Perp^d F \mid \{A, H\}$ hold in the graphs $G, G^{\setminus \{C,E\}},G^{\setminus \{C,E,D\}}$? Explain your answers without using Lemma \ref*{lem:stability_separation_marginalization}, and then check your answer with Lemma \ref*{lem:stability_separation_marginalization}.

    % \emph{Hint: The graphs $G^{\setminus \{C,E\}}$ and $G^{\setminus \{C,E,D\}}$ are displayed in the answer sheet of the exercises of week 4.}
    % \end{enumerate}

    \begin{ungradedsolution}
        % \begin{enumerate}
        %     \item
        \begin{enumerate}
            \item Yes
            \item No, the walk $C\rightarrow D \leftarrow E$ is active because the collider $D$ is in the conditioning set.
            \item No, the walk $B\leftarrow C \to D$ is active as the confounder $C$ is not in the conditioning set.
            \item Yes, the confounder $C$ is in the conditioning set.
        \end{enumerate}
        %     \item Yes. In $G$, the walk $B\leftarrow C \rightarrow D\leftarrow E\rightarrow F$ is blocked by a collider $D$ which is not in the conditioning set. In $G^{\setminus \{C,E\}}$, the walks $B\leftrightarrow D\leftrightarrow F$ and $B\leftrightarrow D\leftarrow H \rightarrow F$ are blocked by collider $D$ as well, and in $G^{\setminus \{C,E,D\}}$ there are no walks between $B$ and $F$. So in all three graphs the $d$-separation holds, which is in accordance with Lemma \ref*{lem:stability_separation_marginalization}.
        % \end{enumerate}
    \end{ungradedsolution}

    \item %(2 + 1 points)
    Let the graphs $G$ and $G'$ be given by
    \begin{figure}[!h]
        \centering
        \begin{subfigure}{0.4\linewidth}
            \centering
            \begin{tikzpicture}[scale=1, transform shape]
                \node[ndout] (X1) at (0, 0) {$X_1$};
                \node[ndout] (X2) at (2, 0) {$X_2$};
                \node[ndout] (X3) at (1, -1.5) {$X_3$};
                \node[ndout] (X4) at (3, -1.5) {$X_4$};
                \node[ndout] (X5) at (2, -3) {$X_5$};
                \node[ndout] (X6) at (4.7, -1.5) {$X_6$};
                \draw[arout] (X1) to (X2);
                \draw[arout] (X2) to (X4);
                \draw[arout] (X4) to (X3);
                \draw[arout] (X3) to (X2);
                \draw[arout] (X4) to (X6);
                \draw[arlat] (X3) to (X5);
            \end{tikzpicture}
            \caption{$G$}
        \end{subfigure}
        \begin{subfigure}{0.4\linewidth}
            \centering
            \begin{tikzpicture}[scale=1, transform shape]
                \node[ndout] (X1) at (0, 0) {$X_1$};
                \node[ndout] (X2) at (2, 0) {$X_2$};
                \node[ndout] (X3) at (1, -1.5) {$X_3$};
                \node[ndout] (X4) at (3, -1.5) {$X_4$};
                \node[ndout] (X5) at (2, -3) {$X_5$};
                \node[ndout] (X6) at (4.7, -1.5) {$X_6$};
                \draw[arout] (X1) to (X2);
                \draw[arout] (X1) to (X3);
                \draw[arout] (X1) to (X4);
                \draw[arout] (X2) to (X4);
                \draw[arout] (X4) to (X3);
                \draw[arout] (X2) to (X3);
                \draw[arout] (X4) to (X6);
                \draw[arlat] (X3) to (X5);
                \draw[arlat] (X2) to (X5);
                \draw[arlat] (X4) to (X5);
            \end{tikzpicture}
            \caption{$G'$}
        \end{subfigure}
    \end{figure}
    \begin{enumerate}
        \item Verify using to Definition \ref*{def:cdg_acyclification} that $G'$ is an acyclification of the graph $G$.
        \item Do the following $d/\sigma$-separations hold?
        \begin{enumerate}
            \item $\displaystyle X_1\Perp^\sigma_{G} X_6|X_2$
            \item $\displaystyle X_1\Perp^\sigma_{G'} X_6|X_2$
            \item $\displaystyle X_1\Perp^d_{G'} X_6|X_2$
        \end{enumerate}
        Whenever $X_1\nPerp X_6|X_2$, provide a $d/\sigma$-open walk. Finally, compare your answers with Proposition \ref*{prp:sigma_sep_as_d_sep}.
    \end{enumerate}

    \begin{ungradedsolution}
        \begin{enumerate}
            \item -
            \item
            \begin{enumerate}
                \item No, the walk $X_1 \to X_2 \to X_4 \to X_6$ is open because neither end points are in the conditioning set, $X_2$ is in the conditioning set but points to an element within its own strongly connected component, and $X_4$ points to another strongly connected component is not in the conditioning set.
                \item The walk $X_1\to X_4 \to X_6$ is open, so no.
                \item The walk $X_1\to X_4 \to X_6$ is open, so no.
            \end{enumerate}
        \end{enumerate}
    \end{ungradedsolution}

    \item
    \begin{enumerate}
        \item Construct an acyclic directed mixed graph consisting only of three (non-input) nodes $A,B,C$, such that no edge (neither directed nor bidirected) exists between $A$ and $C$ and we have:
        \[
            A \nPerp^d C \quad\text{and}\quad A \nPerp^d C \mid B
        \]
        \item %(1 point)
        For graphs with input nodes, the symmetry axiom
        \[
            A \Perp^d B \mid C \implies B \Perp^d A \mid C
        \]
        does not hold in general. Write down a graph for which this implication does not hold and explain why.
        \item  %(1 point)
        Write down a graph for which the $\sigma$-separations and $d$-separations do not agree and explain why.
    \end{enumerate}

    \begin{gradedsolution}
        \begin{enumerate}
            \item \points{1} For example:
            \[
                \begin{tikzpicture}[scale=1, transform shape]
                    \node[ndout] (A) at (0, 0) {$A$};
                    \node[ndout] (B) at (2, 0) {$B$};
                    \node[ndout] (C) at (4, 0) {$C$};
                    \draw[arout] (A) to (B);
                    \draw[arout] (B) to (C);
                    \draw[arlat, bend left] (B) to (C);
                \end{tikzpicture}
            \]
            \item \points{1} For example:
            \[
                \begin{tikzpicture}[scale=1, transform shape]
                    \node[ndout] (B) at (2, 0) {$B$};
                    \node[ndout] (C) at (4, 0) {$C$};
                    \node[ndout] (A) at (6, 0) {$A$};
                    \node[ndint] (I) at (0, 0) {$I$};
                    \draw[arout] (I) to (B);
                    \draw[arout] (B) to (C);
                    \draw[arout] (C) to (A);
                \end{tikzpicture}
            \]
            Here we have $\displaystyle A \Perp^d B \mid C$, but $\displaystyle B \nPerp^d A \mid C$.
            \item \points{1} For example:
            \[
                \begin{tikzpicture}[scale=1, transform shape]
                    \node[ndout] (A) at (0, 0) {$A$};
                    \node[ndout] (B) at (2, 0) {$B$};
                    \node[ndout] (C) at (4, 0) {$C$};
                    \node[ndout] (D) at (3, -1) {$D$};
                    \draw[arout] (A) to (B);
                    \draw[arout] (B) to (C);
                    \draw[arout] (C) to (D);
                    \draw[arout] (D) to (B);
                \end{tikzpicture}
            \]
            The path $A\to B\to C$ is $d$-blocked by a non-collider $B$ which is in the conditioning set, and the path $C \to D \to B \leftarrow A$ is blocked given $B$ and $D$, so $\displaystyle A \Perp^d C \mid B, D$.
            % TODO next year: this is wrong: the path $A \to B \ot D \ot C$ is open? Instead consider A \Perp D | C, this is d-separated sigma-open.

            On the other hand, the walk $A\rightarrow B\rightarrow C$ is $\sigma$-open, since the non-collider $B$ points to an element $C$ within its own strongly connected component $\Sc(B) = \{B, C, D\}$, and therefore $\displaystyle A \nPerp^\sigma C \mid B, D$.
        \end{enumerate}
    \end{gradedsolution}

\end{enumerate}
\end{document}
